\chapter[Léxico de negação do \emph{baseline} de regras]{Léxico de negação do \emph{baseline} de regras}
\label{ap:lexico}

O \emph{baseline} de regras léxicas (B1), descrito na
Seção~\ref{cap:prototipo}, identifica a relação \texttt{negation\_of} a partir de
um conjunto de marcadores de negação adaptados ao português clínico, na linha do
algoritmo NegEx \cite{chapman_2001_negex}. A detecção é \emph{case-insensitive} e
emprega fronteira de palavra, de modo a evitar correspondências espúrias (por
exemplo, ``sem'' dentro de ``semelhante''). Os marcadores são buscados na janela
entre as duas entidades e em uma janela de contexto de $30$ caracteres à esquerda
da segunda entidade.

\begin{quadro}[htbp]
  \centering
  \caption{Marcadores de negação utilizados no \emph{baseline} B1.}
  \label{quad:lexico}
  \begin{tabular}{ll}
  \toprule
  Marcador & Observação \\
  \midrule
  não       & negação verbal/adverbial geral \\
  sem       & preposição de ausência \\
  nega      & forma verbal (``paciente nega\ldots'') \\
  ausência  & substantivo de ausência \\
  ausente   & adjetivo de ausência \\
  afebril   & adjetivo com negação léxico-morfológica (sem febre) \\
  anictérico & adjetivo com negação léxico-morfológica (sem icterícia) \\
  \bottomrule
  \end{tabular}\\[4pt]
  {\footnotesize Fonte: elaborado pelo autor (\texttt{scripts/baselines/baseline\_rules.py}).}
\end{quadro}

O Quadro~\ref{quad:lexico} lista os marcadores. Os termos ``afebril'' e
``anictérico'' são adjetivos que embutem negação morfologicamente e que, no
SemClinBr, aparecem anotados como entidades de negação, o que justifica seu uso
como gatilho. As limitações dessa lista --- ausência de tratamento de negação de
escopo amplo e de negação estendida por conjunção --- são discutidas no
Capítulo~\ref{cap:discussao}.
